% Part: first-order-logic
% Chapter: natural-deduction
% Section: soundness

\documentclass[../../../include/open-logic-section]{subfiles}

\begin{document}

\iftag{FOL}
      {\olfileid[fr]{fol}{ntd}{sou}}
      {\olfileid[fr]{pl}{ntd}{sou}}
      
\olsection{Correction}

\begin{explain}
Un système de !!{derivation}, tel que la déduction naturelle, est
\emph{correct} s'il ne permet de !!{derive} que ce qui résulte
effectivement des hypothèses. La correction garantit ainsi la
fiabilité du système. Selon la propriété de théorie de la
démonstration considérée, nous voulons par exemple savoir que :
\begin{enumerate}
\item tout !!{sentence} !!{derivable} sans hypothèses est
  \iftag{FOL}{valide}{une tautologie};
\item tout !!{sentence} !!{derivable} à partir d'autres est
  aussi une conséquence de ceux-ci;
\item tout ensemble incohérent d'!!{sentence}s est insatisfaisable.
\end{enumerate}
Ces propriétés sont essentielles pour un système de !!{derivation}.
Si l'une d'elles était en défaut, le système permettrait de dériver
trop d'!!{sentence}s. Il est donc primordial d'en établir la
correction.
\end{explain}

\begin{thm}[Correction]
\ollabel{thm:soundness}
Si $!A$ est !!{derivable} à partir d'hypothèses
!!{undischarged}s appartenant à $\Gamma$, alors $\Gamma \Entails !A$.
\end{thm}

\begin{proof}
Soit $\delta$ une !!{derivation} de $!A$. Nous procédons par
récurrence sur le nombre d'inférences de~$\delta$.

Pour le cas de base, supposons ce nombre égal à~$0$. La
!!{derivation}~$\delta$ se réduit alors à un seul !!{sentence}~$!A$,
c'est-à-dire à une hypothèse. Celle-ci est !!{undischarged},
puisqu'une hypothèse ne peut être déchargée que par une inférence
et qu'il n'y en a aucune. Ainsi, toute
\iftag{FOL}{!!{structure}~$\Struct{M}$}{!!{valuation}~$\pAssign{v}$}
qui satisfait toutes les hypothèses !!{undischarged}s de la preuve
satisfait aussi~$!A$.

Passons à l'étape de récurrence. Supposons que $\delta$ comporte~$n$
inférences. Les prémisses de l'inférence la plus basse sont dérivées
par des sous-!!{derivation}s comportant chacune moins de~$n$
inférences. Par hypothèse de récurrence, ces prémisses sont des
conséquences des hypothèses !!{undischarged}s des
sous-!!{derivation}s correspondantes. Il faut montrer que la
conclusion~$!A$ est une conséquence des hypothèses
!!{undischarged}s de la preuve entière.

Distinguons les cas selon la dernière inférence. Commençons par
les inférences à une seule prémisse.

\begin{enumerate}
\item Supposons que la dernière inférence soit \Intro{\lnot}.
  La !!{derivation} a la forme suivante :
\begin{prooftree}
    \AxiomC{$\Gamma, \Discharge{!A}{n}$}
    \RightLabel{$\delta_1$}
    \DeduceC{$\lfalse$}
    \DischargeRule{\Intro{\lnot}}{n}
    \UnaryInfC{$\lnot !A$}
  \end{prooftree}
Par hypothèse de récurrence, $\lfalse$ est une conséquence des
  hypothèses !!{undischarged}s de~$\delta_1$, contenues dans
  $\Gamma \cup \{!A\}$. Considérons
  \iftag{FOL}{!!a{structure}~$\Struct{M}$}{!!a{valuation}~$\pAssign{v}$}.
  Il faut montrer que, si
  $\iftag{FOL}{\Sat{M}}{\pSat{v}}{\Gamma}$, alors
  $\iftag{FOL}{\Sat{M}}{\pSat{v}}{\lnot !A}$.
  Supposons par l'absurde que
  $\iftag{FOL}{\Sat{M}}{\pSat{v}}{\Gamma}$, mais que
  $\iftag{FOL}{\Sat/{M}}{\pSat/{v}}{\lnot !A}$, c'est-à-dire que
  $\iftag{FOL}{\Sat{M}}{\pSat{v}}{!A}$. Nous aurions alors
  $\iftag{FOL}{\Sat{M}}{\pSat{v}}{\Gamma \cup \{!A\}}$, ce qui
  contredit l'hypothèse de récurrence, puisque $\lfalse$ ne peut
  être satisfait. Donc
  $\iftag{FOL}{\Sat{M}}{\pSat{v}}{\lnot !A}$.

\item La dernière inférence est \Elim{\land}. Il y a deux
  variantes : la prémisse $!A \land !B$ permet de conclure $!A$
  ou $!B$. Dans le premier cas, la !!{derivation}~$\delta$ a la forme :
\begin{prooftree}
    \AxiomC{$\Gamma$}
    \RightLabel{$\delta_1$}
    \DeduceC{$!A \land !B$}
    \RightLabel{\Elim{\land}}
    \UnaryInfC{$!A$}
  \end{prooftree}
Par hypothèse de récurrence, $!A \land !B$ est une conséquence
  des hypothèses !!{undischarged}s~$\Gamma$ de~$\delta_1$.
  Considérons
  \iftag{FOL}{!!a{structure}~$\Struct{M}$}{!!a{valuation}~$\pAssign{v}$}.
  Il faut montrer que, si
  $\iftag{FOL}{\Sat{M}}{\pSat{v}}{\Gamma}$, alors
  $\iftag{FOL}{\Sat{M}}{\pSat{v}}{!A}$.
  Supposons $\iftag{FOL}{\Sat{M}}{\pSat{v}}{\Gamma}$.
  Par hypothèse de récurrence, $\Gamma \Entails !A \land !B$;
  nous avons donc $\iftag{FOL}{\Sat{M}}{\pSat{v}}{!A \land !B}$.
  Par définition, cette dernière condition équivaut à
  $\iftag{FOL}{\Sat{M}}{\pSat{v}}{!A}$ et
  $\iftag{FOL}{\Sat{M}}{\pSat{v}}{!B}$.
  Le cas où l'on conclut $!B$ à partir de $!A \land !B$ se traite
  de la même façon.

\item La dernière inférence est \Intro{\lor}. On peut conclure
  $!A \lor !B$ à partir de~$!A$ ou à partir de~$!B$.
  Dans le premier cas, la !!{derivation} a la forme suivante :
\begin{prooftree}
    \AxiomC{$\Gamma$}
    \RightLabel{$\delta_1$}
    \DeduceC{$!A$}
    \RightLabel{\Intro{\lor}}
    \UnaryInfC{$!A \lor !B$}
  \end{prooftree}
Par hypothèse de récurrence, $!A$ est une conséquence des
  hypothèses !!{undischarged}s~$\Gamma$ de~$\delta_1$.
  Considérons
  \iftag{FOL}{!!a{structure}~$\Struct{M}$}{!!a{valuation}~$\pAssign{v}$}.
  Il faut montrer que, si
  $\iftag{FOL}{\Sat{M}}{\pSat{v}}{\Gamma}$, alors
  $\iftag{FOL}{\Sat{M}}{\pSat{v}}{!A \lor !B}$.
  Supposons $\iftag{FOL}{\Sat{M}}{\pSat{v}}{\Gamma}$.
  Comme $\Gamma \Entails !A$ par hypothèse de récurrence,
  $\iftag{FOL}{\Sat{M}}{\pSat{v}}{!A}$, d'où
  $\iftag{FOL}{\Sat{M}}{\pSat{v}}{!A \lor !B}$.
  Le cas où l'on conclut $!A \lor !B$ à partir de~$!B$ est analogue.

\item La dernière inférence est \Intro{\lif}. L'!!{sentence}
  $!A \lif !B$ est dérivé d'une sous-preuve de conclusion~$!B$
  comportant éventuellement des hypothèses~$!A$ à décharger :
\begin{prooftree}
    \AxiomC{$\Gamma, \Discharge{!A}{n}$}
    \RightLabel{$\delta_1$}
    \DeduceC{$!B$}
    \DischargeRule{\Intro{\lif}}{n}
    \UnaryInfC{$!A \lif !B$}
  \end{prooftree}
Par hypothèse de récurrence, $!B$ est une conséquence des
  hypothèses !!{undischarged}s de~$\delta_1$; ainsi,
  $\Gamma \cup \{!A\} \Entails !B$.
  Considérons
  \iftag{FOL}{!!a{structure}~$\Struct{M}$}{!!a{valuation}~$\pAssign{v}$}.
  Les occurrences indiquées de~$!A$ étant déchargées par la
  dernière inférence, les hypothèses !!{undischarged}s de~$\delta$
  appartiennent à~$\Gamma$. Il faut donc montrer
  $\Gamma \Entails !A \lif !B$.
  Supposons par l'absurde qu'il existe
  \iftag{FOL}{!!a{structure}~$\Struct{M}$}{!!a{valuation}~$\pAssign{v}$}
  telle que $\iftag{FOL}{\Sat{M}}{\pSat{v}}{\Gamma}$ et
  $\iftag{FOL}{\Sat/{M}}{\pSat/{v}}{!A \lif !B}$.
  Alors $\iftag{FOL}{\Sat{M}}{\pSat{v}}{!A}$ et
  $\iftag{FOL}{\Sat/{M}}{\pSat/{v}}{!B}$.
  Mais $!B$ est une conséquence de $\Gamma \cup \{!A\}$, donc
  $\iftag{FOL}{\Sat{M}}{\pSat{v}}{!B}$, contradiction.
  Par conséquent, $\Gamma \Entails !A \lif !B$.

\item La dernière inférence est \FalseInt. La !!{derivation}~$\delta$
  se termine alors ainsi :
\begin{prooftree}
    \AxiomC{$\Gamma$}
    \RightLabel{$\delta_1$}
    \DeduceC{$\lfalse$}
    \RightLabel{\FalseInt}
    \UnaryInfC{$!A$}
  \end{prooftree}
Par hypothèse de récurrence, $\Gamma \Entails \lfalse$.
  Il faut montrer que $\Gamma \Entails !A$.
  Sinon, il existerait
  \iftag{FOL}{une structure~$\Struct{M}$}{une valuation~$\pAssign{v}$}
  telle que $\iftag{FOL}{\Sat{M}}{\pSat{v}}{\Gamma}$ et
  $\iftag{FOL}{\Sat/{M}}{\pSat/{v}}{!A}$.
  Or $\iftag{FOL}{\Sat/{M}}{\pSat/{v}}{\lfalse}$ dans tous les
  cas; on aurait donc $\Gamma \Entails/ \lfalse$, contrairement
  à l'hypothèse de récurrence.

\item La dernière inférence est \FalseCl : exercice.

\iftag{FOL}{
\item La dernière inférence est \Intro{\lforall}. La
  !!{derivation}~$\delta$ a alors la forme suivante :
\begin{prooftree}
    \AxiomC{$\Gamma$}
    \RightLabel{$\delta_1$}
    \DeduceC{$!A(a)$}
    \RightLabel{\Intro{\lforall}}
    \UnaryInfC{$\lforall[x][!A(x)]$}
  \end{prooftree}
Par hypothèse de récurrence, la prémisse $!A(a)$ est une
  conséquence des hypothèses !!{undischarged}s~$\Gamma$.
  Considérons une structure~$\Struct{M}$ telle que
  $\Sat{M}{\Gamma}$. Il faut montrer que
  $\Sat{M}{\lforall[x][!A(x)]}$.
  Puisque $\lforall[x][!A(x)]$ est un !!{sentence}, cela revient
  à montrer que $\Sat{M}{!A(x)}[s]$ pour toute affectation de
  variables~$s$ (\olref[syn][ass]{prop:sat-quant}).
  L'ensemble $\Gamma$ ne contient que des énoncés; nous avons donc
  $\Sat{M}{!B}[s]$ pour tout $!B \in \Gamma$, par la
  \olref[syn][sat]{defn:satisfaction}.
  Soit $\Struct{M'}$ la structure qui ne diffère de $\Struct{M}$
  que par l'interprétation de $a$, fixée à $\Assign{a}{M'} = s(x)$.
  Puisque $a$ ne figure pas dans~$\Gamma$, on a
  $\Sat{M'}{\Gamma}$ par le
  \olref[syn][ext]{cor:extensionality-sent}.
  Comme $\Gamma \Entails !A(a)$, il vient $\Sat{M'}{!A(a)}$.
  L'!!{sentence} $!A(a)$ étant clos, nous avons aussi
  $\Sat{M'}{!A(a)}[s]$, par la
  \olref[syn][ass]{prop:sentence-sat-true}.
  Par la \olref[syn][ext]{prop:ext-formulas},
  $\Sat{M'}{!A(x)}[s]$ si et seulement si $\Sat{M'}{!A(a)}$
  (rappelons que $!A(a)$ abrège $\Subst{!A(x)}{a}{x}$).
  Donc $\Sat{M'}{!A(x)}[s]$.
  Puisque $a$ ne figure pas dans~$!A(x)$,
  la \olref[syn][ext]{prop:extensionality} donne $\Sat{M}{!A(x)}[s]$.
  L'affectation~$s$ étant arbitraire, on obtient
  $\Sat{M}{\lforall[x][!A(x)]}$.

\item La dernière inférence est \Intro{\lexists} : exercice.

\item La dernière inférence est \Elim{\lforall} : exercice.
}{}
\end{enumerate}

Considérons maintenant les inférences à plusieurs prémisses :
\Elim{\lor}, \Intro{\land}, \Elim{\lnot},
\iftag{FOL}{\Elim{\lif} et \Elim{\lexists}}{\Elim{\lif}}.
\begin{enumerate}

\item La dernière inférence est \Intro{\land}. On dérive
  $!A \land !B$ des prémisses $!A$ et $!B$, et $\delta$ a la forme :
\begin{prooftree}
    \AxiomC{$\Gamma_1$}
    \RightLabel{$\delta_1$}
    \DeduceC{$!A$}
    \AxiomC{$\Gamma_2$}
    \RightLabel{$\delta_2$}
    \DeduceC{$!B$}
    \RightLabel{\Intro{\land}}
    \BinaryInfC{$!A \land !B$}
  \end{prooftree}
Par hypothèse de récurrence, $!A$ est une conséquence des
  hypothèses !!{undischarged}s~$\Gamma_1$ de~$\delta_1$, et
  $!B$ une conséquence des hypothèses
  !!{undischarged}s~$\Gamma_2$ de~$\delta_2$.
  Les hypothèses !!{undischarged}s de~$\delta$ sont dans
  $\Gamma_1 \cup \Gamma_2$; il faut donc montrer que
  $\Gamma_1 \cup \Gamma_2 \Entails !A \land !B$.
  Considérons
  \iftag{FOL}{!!a{structure}~$\Struct{M}$}{!!a{valuation}~$\pAssign{v}$}
  telle que $\iftag{FOL}{\Sat{M}}{\pSat{v}}{\Gamma_1 \cup \Gamma_2}$.
  Comme $\iftag{FOL}{\Sat{M}}{\pSat{v}}{\Gamma_1}$ et
  $\Gamma_1 \Entails !A$, nous avons
  $\iftag{FOL}{\Sat{M}}{\pSat{v}}{!A}$.
  De même, $\iftag{FOL}{\Sat{M}}{\pSat{v}}{\Gamma_2}$ et
  $\Gamma_2 \Entails !B$ donnent
  $\iftag{FOL}{\Sat{M}}{\pSat{v}}{!B}$.
  Ainsi, $\iftag{FOL}{\Sat{M}}{\pSat{v}}{!A \land !B}$.

\item La dernière inférence est \Elim{\lor} : exercice.

\item La dernière inférence est \Elim{\lif}. On dérive $!B$
  des prémisses $!A \lif !B$ et~$!A$. La !!{derivation}~$\delta$
  a la forme suivante :
\begin{prooftree}
    \AxiomC{$\Gamma_1$}
    \RightLabel{$\delta_1$}
    \DeduceC{$!A \lif !B$}
    \AxiomC{$\Gamma_2$}
    \RightLabel{$\delta_2$}
    \DeduceC{$!A$}
    \RightLabel{\Elim{\lif}}
    \BinaryInfC{$!B$}
  \end{prooftree}
Par hypothèse de récurrence, $!A \lif !B$ est une conséquence
  des hypothèses !!{undischarged}s~$\Gamma_1$ de~$\delta_1$, et
  $!A$ une conséquence des hypothèses
  !!{undischarged}s~$\Gamma_2$ de~$\delta_2$.
  Considérons
  \iftag{FOL}{!!a{structure}~$\Struct{M}$}{!!a{valuation}~$\pAssign{v}$}.
  Il faut montrer que, si
  $\iftag{FOL}{\Sat{M}}{\pSat{v}}{\Gamma_1 \cup \Gamma_2}$, alors
  $\iftag{FOL}{\Sat{M}}{\pSat{v}}{!B}$.
  Supposons $\iftag{FOL}{\Sat{M}}{\pSat{v}}{\Gamma_1 \cup \Gamma_2}$.
  Comme $\Gamma_1 \Entails !A \lif !B$, nous avons
  $\iftag{FOL}{\Sat{M}}{\pSat{v}}{!A \lif !B}$.
  Comme $\Gamma_2 \Entails !A$, nous avons
  $\iftag{FOL}{\Sat{M}}{\pSat{v}}{!A}$.
  Il en résulte $\iftag{FOL}{\Sat{M}}{\pSat{v}}{!B}$.
  En effet, si $\iftag{FOL}{\Sat/{M}}{\pSat/{v}}{!B}$, alors
  $\iftag{FOL}{\Sat{M}}{\pSat{v}}{!A}$ entraînerait
  $\iftag{FOL}{\Sat/{M}}{\pSat/{v}}{!A \lif !B}$, en contradiction
  avec $\iftag{FOL}{\Sat{M}}{\pSat{v}}{!A \lif !B}$.

\item La dernière inférence est \Elim{\lnot} : exercice.

\tagitem{FOL}{La dernière inférence est \Elim{\lexists} : exercice.}{}
\end{enumerate}
\end{proof}

\tagprob{FOL}
\begin{prob}
Compléter la démonstration du \olref[fol][ntd][sou]{thm:soundness}.
\end{prob}
\tagendprob

\tagprob{notFOL}
\begin{prob}
Compléter la démonstration du \olref[pl][ntd][sou]{thm:soundness}.
\end{prob}
\tagendprob

\begin{cor}
\ollabel{cor:weak-soundness}
Si $\Proves !A$, alors $!A$ est \iftag{FOL}{valide}{une tautologie}.
\end{cor}

\begin{cor}
\ollabel{cor:consistency-soundness}
Si $\Gamma$ est satisfaisable, alors $\Gamma$ est cohérent.
\end{cor}

\begin{proof}
Raisonnons par contraposition. Supposons $\Gamma$ incohérent.
Alors $\Gamma \Proves \lfalse$ : il existe une !!{derivation}
de $\lfalse$ dont les hypothèses !!{undischarged}s appartiennent
à~$\Gamma$. Par le \olref{thm:soundness}, toute
\iftag{FOL}{!!{structure}~$\Struct{M}$}{!!{valuation}~$\pAssign{v}$}
qui satisfait $\Gamma$ satisfait aussi $\lfalse$.
Or $\iftag{FOL}{\Sat/{M}}{\pSat/{v}}{\lfalse}$ pour toute
\iftag{FOL}{structure~$\Struct{M}$}{valuation~$\pAssign{v}$}.
Aucune \iftag{FOL}{structure~$\Struct{M}$}{valuation~$\pAssign{v}$}
ne satisfait donc $\Gamma$; cet ensemble est insatisfaisable.
\end{proof}

\begin{editorial}
La première propriété de l'introduction concerne explicitement
la dérivabilité sans hypothèses. Dans le cas de l'élimination de
la conjonction, « structure » est adapté en « valuation » pour la
branche propositionnelle. Le nom de la règle d'élimination
universelle utilise la même commande logique que les autres règles.
L'élimination de la négation, traitée ensuite, est ajoutée à la
liste initiale des règles à plusieurs prémisses. Les cas proposés
en exercice restent à compléter par le lecteur.
\end{editorial}

\end{document}
